\documentclass[../../main.tex]{subfiles}

\begin{document}

\chapter{Kriptografi Modern Simetris: Studi Algoritma Block Cipher Lainnya}
\label{chap:block-cipher-lainnya}
% Running head only: the full chapter title is wider than the header line
% (90pt overfull), so the header uses this shorter form.
\chaptermark{Kriptografi Modern Simetris: Block Cipher Lainnya}

\begin{learningoutcome}
Setelah mempelajari bab ini, mahasiswa diharapkan mampu:
\begin{itemize}
    \item Menganalisis struktur dan mekanisme kerja algoritma GOST (Magma) serta membandingkannya dengan DES (Sub-CPMK 1.2).
    \item Menjelaskan fleksibilitas parameter pada algoritma RC5 dan RC6 (Sub-CPMK 1.2).
    \item Memahami karakteristik algoritma \textit{block cipher} lainnya seperti Blowfish, Serpent, Twofish, dan MARS (Sub-CPMK 1.2).
    \item Mengevaluasi perbedaan antara algoritma berbasis Jaringan Feistel dan algoritma berbasis Substitusi-Permutasi (Sub-CPMK 1.2).
\end{itemize}
\end{learningoutcome}

\subfile{section-01}
\subfile{section-02}
\subfile{section-03}
\section{Contoh Terhitung}

\begin{example}[Satu putaran RC5 dengan $w = 16$]
Tinjau RC5 berukuran kata $w = 16$ bit (blok 32 bit) dengan $r = 2$ putaran. Blok plainteks dipecah menjadi $A_0 = \code{0x1234}$ dan $B_0 = \code{0x2345}$, dengan kunci putaran
\[
\begin{aligned}
S_0 &= \code{0x1A2B}, & S_1 &= \code{0x0C0D}, & S_2 &= \code{0x3344},\\
S_3 &= \code{0x5566}, & S_4 &= \code{0x7788}, & S_5 &= \code{0x1122}.
\end{aligned}
\]
Seluruh penjumlahan dilakukan modulo $2^{16}$, dan $\mathrm{ROTL}_{16}(x,r)$ menyatakan rotasi sirkuler kiri sejauh $r \bmod 16$ bit.

\textbf{Langkah 1: penyaringan awal.}
\[ A = A_0 + S_0 = \code{0x1234} + \code{0x1A2B} = \code{0x2C5F} = 11359, \]
\[ B = B_0 + S_1 = \code{0x2345} + \code{0x0C0D} = \code{0x2F52} = 12114. \]

\textbf{Langkah 2: putaran $i = 1$.}
\[ A \leftarrow \mathrm{ROTL}_{16}(A \oplus B,\; B) + S_2. \]
\[ A \oplus B = \code{0x2C5F} \oplus \code{0x2F52} = \code{0x030D} = 781,\qquad B \bmod 16 = 2, \]
\[ \mathrm{ROTL}_{16}(\code{0x030D}, 2) = \code{0x0C34} = 3124,\qquad A = 3124 + 13124 = \code{0x3F78} = 16248. \]
\[ B \leftarrow \mathrm{ROTL}_{16}(B \oplus A,\; A) + S_3. \]
\[ B \oplus A = \code{0x2F52} \oplus \code{0x3F78} = \code{0x102A} = 4138,\qquad A \bmod 16 = 8, \]
\[ \mathrm{ROTL}_{16}(\code{0x102A}, 8) = \code{0x2A10} = 10768,\qquad B = 10768 + 21862 = \code{0x7F76} = 32630. \]

\textbf{Langkah 3: putaran $i = 2$.}
\[ A \leftarrow \mathrm{ROTL}_{16}(\code{0x3F78} \oplus \code{0x7F76},\; \code{0x7F76}) + S_4, \]
\[ \code{0x3F78} \oplus \code{0x7F76} = \code{0x400E} = 16398,\qquad \mathrm{ROTL}_{16}(\code{0x400E}, 6) = \code{0x0390} = 912, \]
\[ A = 912 + 30600 = \code{0x7B18} = 31512. \]
\[ B \leftarrow \mathrm{ROTL}_{16}(\code{0x7F76} \oplus \code{0x7B18},\; \code{0x7B18}) + S_5, \]
\[ \code{0x7F76} \oplus \code{0x7B18} = \code{0x046E} = 1134,\qquad \mathrm{ROTL}_{16}(\code{0x046E}, 8) = \code{0x6E04} = 28164, \]
\[ B = 28164 + 4386 = \code{0x7F26} = 32550. \]

Cipherteks keluarannya adalah $(A,B) = (\code{0x7B18},\code{0x7F26}) = (31512, 32550)$. Kombinasi XOR, rotasi bergantung-data, dan penjumlahan dengan konstanta membuat satu bit plainteks menyebar ke banyak posisi bit cipherteks---inilah sumber difusi RC5.
\end{example}

\begin{example}[Operasi aritmetika pada fungsi $f$ GOST]
Fungsi $f$ GOST menerapkan penjumlahan modulo $2^{32}$ dengan kunci putaran, substitusi delapan kotak-S ($4 \times 4$ bit), lalu pergeseran sirkuler kiri 11 bit.
\begin{enumerate}
    \item \textbf{Penjumlahan modulo $2^{32}$.} Untuk $R_{i-1} = \code{0x0A0B0C0D}$ dan $K_i = \code{0x01020304}$,
    \[ (R_{i-1} + K_i) \bmod 2^{32} = \code{0x0A0B0C0D} + \code{0x01020304} = \code{0x0B0D0F11} = 185405201. \]
    \item \textbf{Pergeseran sirkuler kiri 11 bit.} Misalkan luaran substitusi kedelapan kotak-S adalah $\code{0x80000001}$ (bit ke-31 dan bit ke-0 bernilai 1). Bit ke-31 bergeser ke posisi $(31 + 11) \bmod 32 = 10$, sedangkan bit ke-0 bergeser ke posisi 11, sehingga
    \[ \mathrm{ROTL}_{32}(\code{0x80000001}, 11) = \code{0x00000C00} = 3072. \]
\end{enumerate}
Hasil ini menjelaskan mengapa GOST lebih tahan \textit{brute-force} daripada DES: kunci GOST 256 bit ($2^{256} \approx 1{,}16 \times 10^{77}$) versus DES 56 bit ($2^{56} \approx 7{,}21 \times 10^{16}$), yaitu sekitar $2^{200} \approx 1{,}61 \times 10^{60}$ kali lebih banyak; ditambah 32 putaran (DES 16) dan kotak-S $4\times4$ (DES $6\times4$), kriptanalisis GOST jauh lebih berat.
\end{example}


\begin{obeactivity}
\textbf{Aktivitas 9.1: Perbandingan Struktur Cipher}
1. Buatlah tabel perbandingan antara DES, AES, dan GOST berdasarkan: ukuran blok, panjang kunci, jumlah putaran, dan struktur (Feistel vs Substitusi-Permutasi).
2. Analisis mengapa algoritma dengan jumlah putaran lebih banyak (seperti GOST atau Serpent) cenderung lebih aman tetapi lebih lambat.
3. Diskusikan keuntungan fleksibilitas parameter pada RC5 dibandingkan dengan algoritma yang parameternya tetap.
\end{obeactivity}


\begin{obereflection}
\textbf{Latihan 9.1}
1. Mengapa penggunaan operasi perkalian pada RC6 memberikan keuntungan keamanan dibandingkan hanya menggunakan XOR dan rotasi seperti pada RC5?
2. Jelaskan mengapa GOST dianggap lebih sulit diserang dengan \textit{brute force} dibandingkan DES. Batasi jawaban Anda pada serangan pencarian kunci menyeluruh saja, dan sebutkan secara eksplisit aspek keamanan mana yang tidak tercakup oleh jawaban itu.
3. Apa perbedaan mendasar antara pendekatan desain Blowfish dan Twofish?
\end{obereflection}

\begin{obereflection}
\textbf{Latihan 9.2}
1. Jelaskan mengapa jadwal kunci GOST jauh lebih murah diimplementasikan daripada jadwal
   kunci DES, dan sebutkan konsekuensi keamanan yang timbul dari kesederhanaan itu.
2. Mengapa larik $L$ pada jadwal kunci RC5 dibaca dengan urutan \textit{little-endian}?
   Apa yang terjadi bila sebuah implementasi memakai urutan \textit{big-endian}?
3. Konstanta $P_w$ dan $Q_w$ diturunkan dari bilangan $e$ dan $\phi$. Apakah RC5 akan
   menjadi lebih lemah bila kedua konstanta itu diganti dengan nilai acak lain? Jelaskan
   alasan Anda.
4. Blowfish dan GOST sama-sama berukuran blok 64 bit. Jelaskan bagaimana keterbatasan ini
   memengaruhi cara kedua algoritma itu seharusnya dipakai dalam sistem nyata.
5. Serpent adalah finalis AES yang paling lambat di perangkat lunak, tetapi justru
   dirancang agar dapat dipercepat dengan \textit{bit-slicing}. Jelaskan gagasan teknik
   tersebut dan mengapa pemilihan kotak-S $4 \times 4$ bit menjadi prasyaratnya.
\end{obereflection}

\section{Rangkuman}
\begin{itemize}
    \item GOST (Magma) adalah \textit{block cipher} standar Rusia dengan blok 64 bit, kunci 256 bit, jaringan Feistel 32 putaran, dan 8 kunci internal $K_1,\dots,K_8$ yang dijadwalkan berulang.
    \item Fungsi $f$ GOST adalah penjumlahan modulo $2^{32}$ dengan kunci putaran, substitusi 8 kotak-S ($4 \times 4$ bit), lalu pergeseran sirkuler kiri 11 bit.
    \item Dibanding DES, GOST memiliki panjang kunci yang jauh lebih besar (256 vs 56 bit), jumlah putaran yang lebih banyak (32 vs 16), dan kotak-S yang lebih sederhana ($4 \times 4$ vs $6 \times 4$). Keunggulan itu berlaku pada pencarian kunci menyeluruh dan pada efisiensi implementasi, \emph{bukan} pada ketahanan terhadap kriptanalisis struktural.
    \item Kriptanalisis atas GOST penuh 32 putaran telah menurunkan kompleksitas serangan dari $2^{256}$ menjadi sekitar $2^{178}$. Karena setiap serangan itu menuntut data paling sedikit $2^{32}$ blok --- tepat pada batas rekeying akibat paradoks ulang tahun --- GOST hanya dapat diserang bila kunci dipakai melampaui batas tersebut.
    \item Setiap \textit{block cipher} berukuran blok $n$ bit hanya boleh dipakai untuk mengenkripsi sekitar $2^{n/2}$ blok dengan satu kunci. Untuk blok 64 bit, batas itu adalah $2^{32}$ blok atau sekitar 32 GiB, dan berlaku bagi GOST maupun Blowfish.
    \item RC5 (Rivest, 1994) memiliki parameter fleksibel: ukuran kata $w \in \{16,32,64\}$, jumlah putaran $r \in \{0,\dots,255\}$, dan panjang kunci $b \in \{0,\dots,255\}$ byte.
    \item Jadwal kunci RC5 menyalin kunci ke larik $L$ secara \textit{little-endian}, menginisialisasi $S[0] = P_w$ dan $S[i] = S[i-1] + Q_w$, lalu mencampur $L$ dan $S$ sebanyak $n = 3 \cdot \max(t,c)$ langkah dengan $t = 2r+2$. Konstanta $P_w$ dan $Q_w$ diturunkan dari bilangan $e$ dan $\phi$ sehingga tidak dapat disembunyikan pintu belakang di dalamnya.
    \item Satu putaran RC5 adalah $A \leftarrow \mathrm{ROTL}(A \oplus B, B) + S_{2i}$ dan $B \leftarrow \mathrm{ROTL}(B \oplus A, A) + S_{2i+1}$, memakai $2r+2$ kunci putaran dengan operasi XOR, penjumlahan modulo $2^w$, dan rotasi yang besarnya ditentukan oleh data itu sendiri.
    \item RC6 memperluas RC5 ke blok 128 bit dengan kunci hingga 2040 bit dan menambahkan perkalian modulo $2^w$ untuk mempercepat difusi.
    \item Blowfish memakai larik $P$ berisi 18 kata dan empat kotak-S berisi 256 entri 32 bit yang seluruhnya dibangkitkan dari digit $\pi$ pada saat kunci dipasang. Twofish menggantinya dengan empat kotak-S $8 \times 8$ yang bergantung pada kunci, sebuah matriks MDS, dan transformasi PHT.
    \item Serpent adalah jaringan substitusi--permutasi murni dengan 32 putaran dan delapan kotak-S $4 \times 4$ bit; MARS adalah Feistel tak seimbang yang membagi 32 putarannya menjadi delapan putaran \textit{forward mixing}, enam belas putaran inti, dan delapan putaran \textit{backward mixing}.
\end{itemize}


\begin{obeassessment}
\textbf{Kuis Bab 9}
1. Jelaskan mekanisme kerja fungsi $f$ pada algoritma GOST.
2. Apa yang dimaksud dengan parameter $w, r, b$ pada RC5 dan bagaimana pengaruhnya terhadap performa cipher?
3. Bandingkan antara Serpent dan AES dari sisi jumlah putaran dan filosofi desainnya.
\end{obeassessment}
\begin{competencychecklist}
\begin{tabularx}{\linewidth}{|X|c|}
\hline
Indikator Kompetensi & Tercapai \\
\hline
Saya dapat menjelaskan mekanisme kerja GOST dan perbedaannya dengan DES & [ ] \\
\hline
Saya dapat menguraikan parameter fleksibel pada RC5 dan RC6 & [ ] \\
\hline
Saya dapat membedakan karakteristik Blowfish, Twofish, Serpent, dan MARS & [ ] \\
\hline
Saya dapat menganalisis trade-off antara jumlah putaran dan efisiensi komputasi & [ ] \\
\hline
Saya dapat membedakan struktur Jaringan Feistel dan Substitusi-Permutasi & [ ] \\
\hline
\end{tabularx}
\end{competencychecklist}


\end{document}
